% =====================================================================
\chapter{Identificación de sistemas y gemelos digitales}
\label{tema:2}
% =====================================================================

El MPC necesita un modelo para predecir y el aprendizaje por refuerzo necesita un simulador en el que equivocarse sin coste. En planta casi nunca se dispone de un modelo físico calibrado: hay que construirlo con datos. Este tema formaliza la \emph{identificación de sistemas} siguiendo a Ljung~\cite{ljung}, texto del curso 6.435 del MIT~\cite{mit6435}, extiende los modelos lineales a los no lineales con redes neuronales como en Narendra y Parthasarathy~\cite{narendra} y define el \emph{gemelo digital} según la clasificación de Kritzinger \emph{et al.}~\cite{kritzinger}. El mensaje central es metodológico: un modelo para control se valida simulando, no prediciendo un paso.

\section{El problema de identificación}

Se dispone de un registro $Z^{N}=\{(u_k,y_k)\}_{k=0}^{N}$ de entradas y salidas. Se elige una \emph{estructura de modelo}, una familia de predictores $\hat y_{k+1}(\theta)$ que dependen de parámetros $\theta$ y de los datos hasta el instante $k$, y se estima $\theta$ minimizando el error de predicción~\cite{ljung}:
\begin{equation}\label{eq:pem}
\hat\theta_N=\arg\min_\theta\ V_N(\theta),\qquad V_N(\theta)=\frac1N\sum_{k=0}^{N-1}\bigl(y_{k+1}-\hat y_{k+1}(\theta)\bigr)^{2}.
\end{equation}

\begin{definicion}[Modelo ARX]\label{def:arx}
Un modelo ARX$(n_a,n_b)$, \emph{autorregresivo con entradas exógenas}, es
\[
y_{k+1}=\sum_{i=1}^{n_a}a_iy_{k+1-i}+\sum_{j=1}^{n_b}b_ju_{k+1-j}+\varepsilon_{k+1},
\]
con $\varepsilon$ ruido blanco. Con el \emph{vector de regresores} $\varphi_k=(y_k,\dots,y_{k+1-n_a},u_k,\dots,u_{k+1-n_b})$ y $\theta=(a,b)$, se escribe $y_{k+1}=\varphi_k\T\theta+\varepsilon_{k+1}$: es una regresión lineal.
\end{definicion}

\begin{proposicion}[Estimador de mínimos cuadrados del ARX]\label{prop:ls}
Sea $\Phi$ la matriz cuyas filas son los $\varphi_k\T$ e $Y$ el vector de los $y_{k+1}$. Si $\Phi\T\Phi$ es invertible, el minimizador de \eqref{eq:pem} es $\hat\theta_N=(\Phi\T\Phi)^{-1}\Phi\T Y$. Si además los datos proceden de un ARX con parámetro $\theta_0$, entonces
\[
\hat\theta_N-\theta_0=\Bigl(\frac1N\sum_k\varphi_k\varphi_k\T\Bigr)^{-1}\frac1N\sum_k\varphi_k\varepsilon_{k+1},
\]
y $\hat\theta_N\to\theta_0$ cuando $\frac1N\sum\varphi_k\varphi_k\T$ converge a una matriz invertible y $\frac1N\sum\varphi_k\varepsilon_{k+1}\to0$.
\end{proposicion}

\begin{proof}
La primera parte son las ecuaciones normales (módulo~8, §2.2). Para la segunda, se sustituye $Y=\Phi\theta_0+\varepsilon$: $\hat\theta_N=(\Phi\T\Phi)^{-1}\Phi\T(\Phi\theta_0+\varepsilon)=\theta_0+(\Phi\T\Phi)^{-1}\Phi\T\varepsilon$, y se multiplica y divide por $N$. El último límite se sigue de la continuidad de la inversa.
\end{proof}

La condición $\frac1N\sum\varphi_k\varepsilon_{k+1}\to0$ se cumple si $\varepsilon$ es blanco, porque $\varphi_k$ solo contiene datos hasta $k$. Si el ruido está coloreado (por ejemplo, ruido de medida sobre una salida autorregresiva), $y_k$ y $\varepsilon_{k+1}$ están correlacionados y el ARX es sesgado; la solución es estimar también un modelo del ruido (ARMAX, error de salida) o usar variables instrumentales~\cite{ljung}. La condición sobre $\frac1N\sum\varphi_k\varphi_k\T$ exige que la entrada sea suficientemente rica.

\begin{definicion}[Excitación persistente]\label{def:pe}
Una señal $u$ es \emph{persistentemente excitante de orden $n$} si la matriz de covarianzas $\bigl(R_u(i-j)\bigr)_{i,j=1}^{n}$, con $R_u(\tau)=\lim_N\frac1N\sum_ku_ku_{k-\tau}$, es definida positiva.
\end{definicion}

Una constante es persistentemente excitante solo de orden~1; una sinusoide, de orden~2; un ruido blanco o una secuencia binaria pseudoaleatoria, de cualquier orden. En planta se usan escalones aleatorios de duración variable, que respetan los límites del actuador y excitan las frecuencias que interesan.

\begin{ejemplo}[Un ARX a mano]\label{ej:arx-mano}
Los datos $u=(1,0,1)$ y $y=(0;\ 0{,}5;\ 0{,}45;\ 0{,}905)$ proceden de $y_{k+1}=0{,}9y_k+0{,}5u_k$. Con $\varphi_k=(y_k,u_k)$:
\[
\Phi\T\Phi=\begin{pmatrix}0{,}4525&0{,}45\\0{,}45&2\end{pmatrix},\qquad \Phi\T Y=\begin{pmatrix}0{,}63225\\1{,}405\end{pmatrix},\qquad \det=0{,}7025,
\]
y $\hat a=(2\cdot0{,}63225-0{,}45\cdot1{,}405)/0{,}7025=0{,}9$, $\hat b=(0{,}4525\cdot1{,}405-0{,}45\cdot0{,}63225)/0{,}7025=0{,}5$. Sin ruido, tres ecuaciones bastan para recuperar los dos parámetros exactos.
\end{ejemplo}

\section{Modelos no lineales: NARX y redes neuronales}

\begin{definicion}[Modelo NARX]\label{def:narx}
Un modelo NARX es $y_{k+1}=g(\varphi_k;\theta)+\varepsilon_{k+1}$ con los mismos regresores que el ARX y una función $g$ no lineal: un polinomio, una red de base radial o un perceptrón multicapa (módulo~9). Narendra y Parthasarathy~\cite{narendra} propusieron en 1990 identificar y controlar sistemas dinámicos con redes de este tipo.
\end{definicion}

Hay tres familias principales, de menos a más estructura física:
\begin{enumerate}
  \item \emph{Caja negra}: $g$ es un perceptrón multicapa sobre los regresores. Es flexible, pero solo es fiable dentro del rango de los datos.
  \item \emph{Modelos de estado neuronales}: un estado latente $h_{k+1}=F(h_k,u_k)$ e $y_k=H(h_k)$ con $F$ y $H$ redes, es decir, una red recurrente (módulo~12, §5.2) entrenada en simulación sobre secuencias, con retropropagación en el tiempo. Captura dinámicas de orden alto sin elegir $n_a$, a costa de un entrenamiento más delicado.
  \item \emph{Caja gris}: la física aporta la estructura y los datos los parámetros, o una red aprende solo el residuo del modelo físico. Si se conoce la forma no lineal (por ejemplo, un producto entre entrada y perturbación), basta con añadir el regresor correspondiente y el modelo sigue siendo lineal en $\theta$.
\end{enumerate}

\begin{ejemplo}[Modelos de Hammerstein]\label{ej:hammerstein}
Si una no linealidad estática actúa sobre la entrada, $y_{k+1}=ay_k+b\,f(u_k)$ con $f(u)=u+cu^{2}$, el modelo es lineal en $(a,\ b,\ bc)$ con regresores $(y_k,\ u_k,\ u_k^{2})$: se estima por mínimos cuadrados como un sistema de dos entradas y se recupera $c=\widehat{bc}/\hat b$. Solo es identificable el producto $b\,f$: por eso se fija a~1 el coeficiente lineal de $f$ (\cref{eje:2-hammerstein}).
\end{ejemplo}

\section{Validación: predicción a un paso y simulación libre}

\begin{definicion}[Dos formas de validar]\label{def:validacion}
Sea $\hat f$ un modelo de un paso, $\hat x_{k+1}=\hat f(x_k,u_k)$. La \emph{predicción a un paso} usa en cada instante el estado medido: $\hat x_{k+1}=\hat f(x_k,u_k)$. La \emph{simulación libre} (validación en lazo abierto del modelo, o \emph{free run}) se realimenta con su propia salida: $\tilde x_{k+1}=\hat f(\tilde x_k,u_k)$, $\tilde x_0=x_0$.
\end{definicion}

El MPC predice $N$ pasos y el agente de RL se entrena durante episodios enteros en el modelo: ambos usan el modelo en simulación libre. La siguiente proposición explica por qué el error a un paso es engañoso.

\begin{proposicion}[Acumulación del error en simulación]\label{prop:acumula}
Sea $x_{k+1}=f(x_k,u_k)$ el sistema real y $\hat f$ un modelo tal que $\abs{\hat f(x,u)-f(x,u)}\le\varepsilon$ para todo $(x,u)$ y $\hat f$ es lipschitziana en $x$ con constante $L$. Entonces el error de simulación libre cumple
\[
\abs{\tilde x_k-x_k}\le\varepsilon\,\frac{1-L^{k}}{1-L}\quad(L\ne1),\qquad\text{y, si }L<1,\quad \abs{\tilde x_k-x_k}\le\frac{\varepsilon}{1-L}.
\]
\end{proposicion}

\begin{proof}
Sea $d_k=\abs{\tilde x_k-x_k}$, con $d_0=0$. Por la desigualdad triangular,
$d_{k+1}=\abs{\hat f(\tilde x_k,u_k)-f(x_k,u_k)}\le\abs{\hat f(\tilde x_k,u_k)-\hat f(x_k,u_k)}+\abs{\hat f(x_k,u_k)-f(x_k,u_k)}\le Ld_k+\varepsilon$.
Por inducción, $d_k\le\varepsilon(1+L+\dots+L^{k-1})=\varepsilon(1-L^{k})/(1-L)$, suma geométrica (módulo~0.6, tema~5), acotada por $\varepsilon/(1-L)$ si $L<1$.
\end{proof}

En procesos lentos $L$ es casi~1. Para la zona, $L=a=0{,}9835$ y $1/(1-L)\approx61$: un error de modelo de $0{,}005$ K por paso, invisible a un paso, puede convertirse en $0{,}30$ K en simulación, y en $0{,}24$ K tras las 96 muestras de un horizonte de 24~h.

\begin{proposicion}[Identificabilidad en lazo cerrado]\label{prop:lazo}
Sea $y_{k+1}=ay_k+bu_k+\varepsilon_{k+1}$ y supóngase que los datos se registran con un controlador proporcional $u_k=-Ky_k$ y sin excitación externa. Entonces $\Phi\T\Phi$ es singular y todos los pares $(a',b')$ con $a'-Kb'=a-Kb$ dan el mismo error de predicción: $a$ y $b$ no son identificables por separado.
\end{proposicion}

\begin{proof}
Las filas de $\Phi$ son $(y_k,\ -Ky_k)$, proporcionales a $(1,-K)$: la segunda columna es $-K$ veces la primera y $\Phi\T\Phi$ tiene rango~1. El predictor vale $a'y_k+b'u_k=(a'-Kb')y_k$, que solo depende de $a'-Kb'$.
\end{proof}

Es la razón por la que no basta con descargar el histórico de un lazo que funciona bien: la entrada es una función de la salida y no aporta información independiente. Las soluciones son añadir excitación (escalones o ruido en la consigna o en la salida del controlador), alternar entre controladores distintos o registrar periodos en manual~\cite{ljung}. El histórico de un PID también es poco informativo por otra razón: el controlador mantiene la planta cerca de un único punto de operación. Lazic \emph{et al.}~\cite{lazic} lo comprobaron en un centro de datos: tres horas de exploración aleatoria y acotada produjeron un modelo mejor que una semana de datos históricos de los PID (\cref{tema:5}).

\begin{observacion}[Partición temporal de los datos]\label{obs:particion}
Las muestras de una serie temporal no son independientes: barajarlas antes de separar entrenamiento y validación (módulo~8, §1.4) filtra información, porque cada muestra de validación tiene vecinas casi idénticas en el entrenamiento. Se separan \emph{bloques de tiempo} (días o semanas completos) y se reserva el último bloque, con condiciones distintas (otra estación del año, otro régimen de producción), como prueba final.
\end{observacion}

\section{El gemelo digital}

\begin{definicion}[Modelo, sombra y gemelo digital]\label{def:gemelo}
Kritzinger \emph{et al.}~\cite{kritzinger} distinguen tres niveles según la automatización del flujo de datos entre el objeto físico y su representación digital: un \emph{modelo digital} intercambia datos con la planta solo de forma manual; una \emph{sombra digital} recibe automáticamente los datos de la planta pero no actúa sobre ella; un \emph{gemelo digital} intercambia datos automáticamente en los dos sentidos.
\end{definicion}

\begin{figure}[h]
\centering
\begin{tikzpicture}[>=Stealth,every node/.style={font=\small},
  caja/.style={draw,rounded corners,minimum height=1cm,minimum width=2.6cm,align=center}]
\node[caja] (planta) {planta\\(sensores, DCS)};
\node[caja,right=2.4cm of planta] (est) {estimación de estado\\y calibración};
\node[caja,right=1.2cm of est] (mod) {modelo\\identificado};
\node[caja,below=1.1cm of est] (mpc) {MPC / política};
\node[caja,below=1.1cm of mod] (usos) {entrenamiento RL,\\análisis «¿y si...?»};
\draw[->] (planta) -- node[above]{medidas} (est);
\draw[->] (est) -- (mod);
\draw[->] (mod) -- (mpc);
\draw[->] (mod) -- (usos);
\draw[->] (mpc) -| node[pos=0.25,below]{consignas} (planta);
\draw[->,dashed] (est.north) to[bend left=20] node[above]{reidentificación} (mod.north);
\end{tikzpicture}
\caption{Flujo de datos de un gemelo digital usado para control.}
\label{fig:gemelo}
\end{figure}

Para control, el gemelo tiene cuatro componentes (\cref{fig:gemelo}): un \emph{modelo} (físico, de datos o híbrido); un \emph{estimador de estado}, porque no todos los estados se miden (el filtro de Kalman o la estimación de horizonte móvil, \cref{tema:3}); un procedimiento de \emph{recalibración} cuando los residuos derivan, y una \emph{interfaz} estándar para que lo usen otros programas: el estándar FMI~\cite{fmi}, que empaqueta un modelo como unidad intercambiable (FMU), o la interfaz de entornos de Gymnasium~\cite{gymnasium}. Para edificios, el simulador EnergyPlus~\cite{crawley} y bibliotecas de Modelica~\cite{modelica} son la referencia física, y marcos como BOPTEST~\cite{boptest} o Sinergym~\cite{sinergym} los envuelven como entornos de prueba para controladores.

\begin{ejemplo}[Caso de estudio: identificación del gemelo de la zona]\label{ej:caso-id}
\emph{Datos.} Se registran 14 días de la zona del \cref{ej:caso-pi} en lazo abierto con escalones aleatorios de potencia (de 1 a 4~h de duración y entre 0 y 3~kW), con ruido de medida de $0{,}1$ K; 1344 muestras. Los días 1 a 10 se usan para estimar y los días 11 a 14 para validar.

\emph{Modelos.} (a) ARX con regresores $(T_k,P_k,\Text{}_{,k},G_k,1)$; (b) NARX con física: añade el regresor $P_k\Text{}_{,k}$, porque por \eqref{eq:rc-d} el coeficiente de $P_k$ es $(1-a)R(2{,}5+0{,}06\,\Text)$; (c) NARX neuronal: un perceptrón de dos capas ocultas de 32 unidades sobre las mismas entradas que el ARX, que predice el incremento de temperatura.

\emph{Resultados} (cuaderno, Parte~3):
\begin{center}\small
\begin{tabular}{@{}lcc@{}}
\toprule
Modelo & RMSE a un paso (K) & RMSE en simulación libre, 4 días (K) \\
\midrule
ARX & $0{,}137$ & $0{,}287$ \\
NARX con física & $0{,}136$ & $0{,}180$ \\
NARX neuronal & $0{,}137$ & $0{,}340$ \\
\bottomrule
\end{tabular}
\end{center}
A un paso los tres son indistinguibles: el error está dominado por el ruido de medida de $T_k$ y $T_{k+1}$, $0{,}1\sqrt{1+a^{2}}=0{,}140$ K. En simulación, el modelo con la estructura correcta es claramente mejor, y recupera los parámetros físicos (por ejemplo, $\hat a=0{,}9827$ frente a $0{,}9835$, es decir, $\hat\tau=14{,}3$ h frente a 15~h). La red neuronal, sin esa estructura y con solo diez días de datos, simula peor que el ARX: más capacidad no compensa menos información. Este modelo con física es el que usa el MPC del \cref{tema:3}.
\end{ejemplo}

\begin{table}[htbp]
\centering\footnotesize
\begin{tabularx}{\textwidth}{@{}>{\raggedright\arraybackslash}p{3.1cm}>{\raggedright\arraybackslash}X>{\raggedright\arraybackslash}p{2.0cm}>{\raggedright\arraybackslash}X@{}}
\toprule
Herramienta & Para qué sirve & Licencia o tipo & Cuándo elegirla \\
\midrule
NumPy / SciPy & mínimos cuadrados, ARX, filtrado & libre (BSD-3) & modelos lineales en los parámetros \\
SysIdentPy & modelos NARMAX, selección de regresores & libre (BSD-3) & NARX polinómicos interpretables \\
scikit-learn & perceptrones, \emph{ridge}, validación & libre (BSD-3) & NARX neuronales pequeños \\
PyTorch & redes recurrentes, modelos de estado neuronales, entrenamiento en simulación & libre (tipo BSD) & modelos neuronales con dinámica latente o híbridos con física \\
MATLAB System Identification Toolbox & ARX, ARMAX, error de salida, subespacios & comercial & entornos ya basados en MATLAB \\
FMI / FMPy, Modelica & encapsular y ejecutar el gemelo & estándar abierto / libre & conectar el gemelo con otras herramientas \\
EnergyPlus~\cite{crawley} & simulación energética detallada de edificios & libre (tipo BSD) & gemelo físico de un edificio; generar datos \\
\bottomrule
\end{tabularx}
\caption{Tecnologías para la identificación y los gemelos digitales.}
\label{tab:tec2}
\end{table}

\begin{observacion}[Errores frecuentes en la práctica]\label{obs:errores2}
(i) \emph{Validar solo a un paso}: todos los modelos parecen buenos (\cref{ej:caso-id}). (ii) \emph{Barajar las muestras} antes de separar entrenamiento y validación (\cref{obs:particion}). (iii) \emph{Identificar con datos en lazo cerrado sin excitación} (\cref{prop:lazo}). (iv) \emph{Extrapolar}: un modelo de caja negra entrenado en invierno no sabe nada del verano, y el optimizador buscará precisamente las regiones donde el modelo se equivoca a su favor; es la versión en control del sobreajuste del modelo sustituto. (v) \emph{Incoherencias de muestreo}: el modelo identificado con $\Dt=15$ min no sirve con $\Dt=5$ min, y los retardos de los registros históricos (marcas de tiempo desalineadas) se convierten en dinámica falsa. (vi) \emph{Fuga de información del futuro}: usar como regresor una medida que en tiempo real aún no existiría, como la temperatura exterior real en lugar de su previsión.
\end{observacion}

\begin{lecturas}
\item L.~Ljung~\cite{ljung}, \emph{System Identification: Theory for the User}, en particular los capítulos sobre estructuras de modelos, métodos de error de predicción y diseño de experimentos.
\item MIT 6.435~\cite{mit6435}, temario y hojas de problemas.
\item K.~S. Narendra y K.~Parthasarathy~\cite{narendra}, «Identification and control of dynamical systems using neural networks».
\item W.~Kritzinger \emph{et al.}~\cite{kritzinger}, «Digital Twin in manufacturing: A categorical literature review and classification».
\end{lecturas}

\section*{Ejercicios del tema 2}
\addcontentsline{toc}{section}{Ejercicios del tema 2}

\begin{ejercicio}{1}{eje:2-unpaso}
Explique con sus palabras por qué un modelo puede tener un error a un paso muy pequeño y, sin embargo, ser inútil para un MPC con horizonte de 24~h. ¿Qué métrica de validación usaría y sobre qué datos?
\end{ejercicio}

\begin{ejercicio}{2}{eje:2-arx}
Con los datos $y=(1;\ 1{,}2;\ 0{,}95;\ 1{,}18)$ y $u=(1,0,1)$, estime por mínimos cuadrados el ARX $y_{k+1}=ay_k+bu_k$: escriba $\Phi\T\Phi$ y $\Phi\T Y$, resuelva el sistema y calcule los residuos. Simule después el modelo estimado desde $y_0=1$ y compare con los datos.
\end{ejercicio}

\begin{ejercicio}{2}{eje:2-escalar}
Sea el proceso $y_k=\theta_0u_k+e_k$, con $e$ ruido blanco gaussiano de varianza $\lambda=1$, independiente de $u$, y el predictor $\hat y_k=\theta u_k$. (a) Obtenga el estimador de mínimos cuadrados $\hat\theta_N$ y su límite cuando $N\to\infty$. (b) Calcule su varianza. (c) Compare la desviación típica con $N=400$ para $u_k=\pm1$ alternante y para $u_k=0{,}1\sin(\omega k)$. ¿Qué enseña sobre la amplitud de la excitación?
\fuente{adaptado de MIT 6.435 \emph{System Identification}, hoja de problemas~5, problema~2~\cite{mit6435}.}
\end{ejercicio}

\begin{ejercicio}{2}{eje:2-hammerstein}
Un actuador tiene una característica no lineal y el proceso sigue $y_{k+1}=ay_k+b\,f(u_k)$ con $f(u)=u+cu^{2}$. (a) Escriba el problema como regresión lineal con dos entradas. (b) Si la estimación da $\hat\theta=(0{,}85;\ 0{,}30;\ -0{,}06)$, obtenga $a$, $b$ y $c$. (c) Explique por qué no se podrían identificar por separado $b$ y los dos coeficientes de $f$ si no se fijara uno de ellos.
\fuente{adaptado de MIT 6.435 \emph{System Identification}, hoja de problemas~5, problema~3 (modelos de tipo Hammerstein)~\cite{mit6435}.}
\end{ejercicio}

\begin{ejercicio}{2}{eje:2-cota}
Un modelo de la zona tiene $L=0{,}98$ y un error de modelo de a lo sumo $\varepsilon=0{,}01$ K por paso. (a) Acote el error de simulación tras 4, 24 y 96 pasos y en el límite. (b) Si el MPC necesita un error menor que $0{,}2$ K en todo el horizonte de 96 pasos, ¿qué $\varepsilon$ máximo puede admitirse? (c) Proponga un margen de seguridad para las restricciones de confort a partir de esa cota.
\end{ejercicio}

\begin{ejercicio}{2}{eje:2-lazo}
Una planta $y_{k+1}=0{,}9y_k+0{,}5u_k+\varepsilon_{k+1}$ funciona con un controlador $u_k=-0{,}4y_k$. (a) ¿Qué modelo de lazo cerrado se observa? (b) Dé tres pares $(a',b')$ que expliquen los datos igual de bien. (c) Si se suma a la salida del controlador una señal $v_k$ blanca, independiente del ruido, de varianza $\sigma_v^{2}>0$, demuestre que $\Phi\T\Phi/N$ converge a una matriz invertible.
\end{ejercicio}

\begin{ejercicio}{3}{eje:2-sesgo}
Sea $x_{k+1}=ax_k+bu_k$ un sistema sin ruido de proceso, con $u$ blanco de varianza $\sigma_u^{2}$ y $\abs a<1$, y sea la medida $y_k=x_k+v_k$ con $v$ blanco de varianza $\sigma_v^{2}$, independiente de $u$. Se estima el ARX $y_{k+1}=\alpha y_k+\beta u_k$. (a) Demuestre que $\hat\alpha$ converge a $a\,\sigma_x^{2}/(\sigma_x^{2}+\sigma_v^{2})$, con $\sigma_x^{2}=b^{2}\sigma_u^{2}/(1-a^{2})$ la varianza estacionaria de $x$. (b) Interprete el resultado (atenuación por error en las variables) y calcule el sesgo de la constante de tiempo para $a=0{,}9$, $b=1$, $\sigma_u=1$ y $\sigma_v=1$.
\end{ejercicio}

\begin{ejercicio}{3}{eje:2-diseno}
Diseñe el protocolo de identificación y validación del gemelo digital de una nave industrial climatizada de la que se tiene un año de datos históricos de su PID y se puede intervenir dos semanas. Indique: experimentos y excitación (con sus límites de seguridad), partición de los datos, estructuras de modelo candidatas, métricas, criterio de aceptación para usarlo en un MPC y procedimiento de recalibración.
\end{ejercicio}
